% GENERATED FILE. Edit canonical lessons, then run npm run generate:books.
% canonical-source-hash: fnv1a64:1cc2e9ab7dc81d3d
% canonical-lessons: TE-C239-nivedika, TE-C239-samstha, TE-C239-citi, TE-C239-pinchanu, TE-C239-oppandam, TE-R239-first-pass-words-from-chapters-235-236, TE-R239-second-pass-words-from-chapters-237-239

\chapter{Report, Organization, Chit, Pension, Agreement}
\label{ch:te-report-organization-chit-pension-agreement}
\hlchaptermodality{239}

\begin{chapteropening}
\textbf{By the end of this chapter:} \emph{I can say a report, an organization, a chit, a pension and an agreement.}

Complete the last lesson of chapter 239: I can say a report, an organization, a chit, a pension and an agreement.
\end{chapteropening}

\section[\texorpdfstring{\te{నివేదిక} (nivēdika)}{నివేదిక (nivēdika)}]{\te{నివేదిక} (nivēdika) — a report}

\label{lesson:TE-C239-nivedika}



\begin{quote}
Before the new one: say the Telugu for science, then the Telugu for physics.
\end{quote}

\subsection*{You'll want to know: \te{నివేదిక}}
\textbf{\te{నివేదిక}} — \emph{nivēdika} — \textquotedblleft{}a report\textquotedblright{}.

\begin{grammarlens}[title={What you've built}]
One more word for this chapter.
\end{grammarlens}

\subsection*{Your turn}
\begin{itemize}
\raggedright
  \item \emph{Say it:} \emph{nivēdika}
  \item \emph{Say it:} \emph{nivēdika}, once more
  \item \emph{Say it:} say \emph{nivēdika}
  \item \emph{Recall:} say the Telugu for science, then the Telugu for physics, then say \emph{nivēdika} again
\end{itemize}

\begin{tcolorbox}[breakable,colback=teal!4,colframe=teal!35!black,title={Before you move on}]
What does \te{నివేదిక} mean? (\textquotedblleft{}A report\textquotedblright{}.) Say it once more.
\end{tcolorbox}
\section[\texorpdfstring{\te{సంస్థ} (saṁstha)}{సంస్థ (saṁstha)}]{\te{సంస్థ} (saṁstha) — an organization}

\label{lesson:TE-C239-samstha}



\begin{quote}
Before the new one: say the Telugu for physics, then the Telugu for a report.
\end{quote}

\subsection*{You'll want to know: \te{సంస్థ}}
\textbf{\te{సంస్థ}} — \emph{saṁstha} — \textquotedblleft{}an organization\textquotedblright{}.

\begin{grammarlens}[title={What you've built}]
One more word for this chapter.
\end{grammarlens}

\subsection*{Your turn}
\begin{itemize}
\raggedright
  \item \emph{Say it:} \emph{saṁstha}
  \item \emph{Say it:} \emph{saṁstha}, once more
  \item \emph{Say it:} say \emph{saṁstha}
  \item \emph{Recall:} say the Telugu for physics, then the Telugu for a report, then say \emph{saṁstha} again
\end{itemize}

\begin{tcolorbox}[breakable,colback=teal!4,colframe=teal!35!black,title={Before you move on}]
What does \te{సంస్థ} mean? (\textquotedblleft{}An organization\textquotedblright{}.) Say it once more.
\end{tcolorbox}
\section[\texorpdfstring{\te{చీటీ} (cīṭī)}{చీటీ (cīṭī)}]{\te{చీటీ} (cīṭī) — a chit, a slip of paper}

\label{lesson:TE-C239-citi}



\begin{quote}
Before the new one: say the Telugu for a report, then the Telugu for an organization.
\end{quote}

\subsection*{You'll want to know: \te{చీటీ}}
\textbf{\te{చీటీ}} — \emph{cīṭī} — \textquotedblleft{}a chit, a slip of paper\textquotedblright{}.

\begin{grammarlens}[title={What you've built}]
One more word for this chapter.
\end{grammarlens}

\subsection*{Your turn}
\begin{itemize}
\raggedright
  \item \emph{Say it:} \emph{cīṭī}
  \item \emph{Say it:} \emph{cīṭī}, once more
  \item \emph{Say it:} say \emph{cīṭī}
  \item \emph{Recall:} say the Telugu for a report, then the Telugu for an organization, then say \emph{cīṭī} again
\end{itemize}

\begin{tcolorbox}[breakable,colback=teal!4,colframe=teal!35!black,title={Before you move on}]
What does \te{చీటీ} mean? (\textquotedblleft{}A chit, a slip of paper\textquotedblright{}.) Say it once more.
\end{tcolorbox}
\section[\texorpdfstring{\te{పింఛను} (piñchanu)}{పింఛను (piñchanu)}]{\te{పింఛను} (piñchanu) — a pension}

\label{lesson:TE-C239-pinchanu}



\begin{quote}
Before the new one: say the Telugu for an organization, then the Telugu for a chit.
\end{quote}

\subsection*{You'll want to know: \te{పింఛను}}
\textbf{\te{పింఛను}} — \emph{piñchanu} — \textquotedblleft{}a pension\textquotedblright{}.

\begin{grammarlens}[title={What you've built}]
One more word for this chapter.
\end{grammarlens}

\subsection*{Your turn}
\begin{itemize}
\raggedright
  \item \emph{Say it:} \emph{piñchanu}
  \item \emph{Say it:} \emph{piñchanu}, once more
  \item \emph{Say it:} say \emph{piñchanu}
  \item \emph{Recall:} say the Telugu for an organization, then the Telugu for a chit, then say \emph{piñchanu} again
\end{itemize}

\begin{tcolorbox}[breakable,colback=teal!4,colframe=teal!35!black,title={Before you move on}]
What does \te{పింఛను} mean? (\textquotedblleft{}A pension\textquotedblright{}.) Say it once more.
\end{tcolorbox}
\section[\texorpdfstring{\te{ఒప్పందం} (oppandaṁ)}{ఒప్పందం (oppandaṁ)}]{\te{ఒప్పందం} (oppandaṁ) — an agreement, a contract}

\label{lesson:TE-C239-oppandam}



\begin{quote}
Before the new one: say the Telugu for a chit, then the Telugu for a pension.
\end{quote}

\subsection*{You'll want to know: \te{ఒప్పందం}}
\textbf{\te{ఒప్పందం}} — \emph{oppandaṁ} — \textquotedblleft{}an agreement, a contract\textquotedblright{}.

\begin{grammarlens}[title={What you've built}]
One more word for this chapter.
\end{grammarlens}

\subsection*{Your turn}
\begin{itemize}
\raggedright
  \item \emph{Say it:} \emph{oppandaṁ}
  \item \emph{Say it:} \emph{oppandaṁ}, once more
  \item \emph{Say it:} say \emph{oppandaṁ}
  \item \emph{Recall:} say the Telugu for a chit, then the Telugu for a pension, then say \emph{oppandaṁ} again
\end{itemize}

\begin{tcolorbox}[breakable,colback=teal!4,colframe=teal!35!black,title={Before you move on}]
What does \te{ఒప్పందం} mean? (\textquotedblleft{}An agreement, a contract\textquotedblright{}.) Say it once more.
\end{tcolorbox}
\section[(dialogue)]{Words from chapters 235-236 — a first pass}

\label{lesson:TE-R239-first-pass-words-from-chapters-235-236}



\begin{quote}
Say the words for a pension and an agreement.
\end{quote}

\subsection*{Your turn}
Say these aloud:

\begin{itemize}
\raggedright
  \item a banknote, loose change, a salary, a profit, a loss
  \item a team, business, a company, the government, marks
\end{itemize}

\begin{tcolorbox}[breakable,colback=teal!4,colframe=teal!35!black,title={Before you move on}]
A banknote? (\textbf{\te{నోటు}}, \emph{nōṭu}.) Loose change? (\textbf{\te{చిల్లర}}, \emph{cillara}.) A salary? (\textbf{\te{జీతం}}, \emph{jītaṁ}.) A profit? (\textbf{\te{లాభం}}, \emph{lābhaṁ}.) A loss? (\textbf{\te{నష్టం}}, \emph{naṣṭaṁ}.) A team? (\textbf{\te{జట్టు}}, \emph{jaṭṭu}.) Business? (\textbf{\te{వ్యాపారం}}, \emph{vyāpāraṁ}.) A company? (\textbf{\te{కంపెనీ}}, \emph{kampenī}.) The government? (\textbf{\te{ప్రభుత్వం}}, \emph{prabhutvaṁ}.) Marks? (\textbf{\te{మార్కులు}}, \emph{mārkulu}.)
\end{tcolorbox}
\section[(dialogue)]{Words from chapters 237-239 — a second pass}

\label{lesson:TE-R239-second-pass-words-from-chapters-237-239}



\begin{quote}
Say the words for a police officer and a driver.
\end{quote}

\subsection*{Your turn}
Say these aloud:

\begin{itemize}
\raggedright
  \item a police officer, a driver, a postman, a bus conductor, a lawyer
  \item a memory, history, geography, science, physics
  \item a report, an organization, a chit, a pension, an agreement
\end{itemize}

\begin{tcolorbox}[breakable,colback=teal!4,colframe=teal!35!black,title={Before you move on}]
A police officer? (\textbf{\te{పోలీసు}}, \emph{pōlīsu}.) A driver? (\textbf{\te{డ్రైవరు}}, \emph{ḍraivaru}.) A postman? (\textbf{\te{పోస్టుమాను}}, \emph{pōsṭumānu}.) A bus conductor? (\textbf{\te{కండక్టరు}}, \emph{kaṇḍakṭaru}.) A lawyer? (\textbf{\te{వకీలు}}, \emph{vakīlu}.) A memory? (\textbf{\te{జ్ఞాపకం}}, \emph{jñāpakaṁ}.) History? (\textbf{\te{చరిత్ర}}, \emph{caritra}.) Geography? (\textbf{\te{భూగోళశాస్త్రం}}, \emph{bhūgōḷaśāstraṁ}.) Science? (\textbf{\te{విజ్ఞానశాస్త్రం}}, \emph{vijñānaśāstraṁ}.) Physics? (\textbf{\te{భౌతికశాస్త్రం}}, \emph{bhautikaśāstraṁ}.) A report? (\textbf{\te{నివేదిక}}, \emph{nivēdika}.) An organization? (\textbf{\te{సంస్థ}}, \emph{saṁstha}.) A chit? (\textbf{\te{చీటీ}}, \emph{cīṭī}.) A pension? (\textbf{\te{పింఛను}}, \emph{piñchanu}.) An agreement? (\textbf{\te{ఒప్పందం}}, \emph{oppandaṁ}.)
\end{tcolorbox}
